\documentclass[12pt]{article}
\usepackage[utf8]{inputenc}
\usepackage[T1]{fontenc}
\usepackage{hyperref}
\title{Semimi Language Resources for K-12 Curriculum Development}
\author{Compiled by Claude Sonnet 5.5 (claude-sonnet-5-5) for Andr\'as Kornai}
\date{}
\begin{document}
\maketitle

\section{Introduction}
Semimi (ISO 639-3 etz; Glottocode semi1264; also called Etna Bay) is a Papuan language of the Mairasi family, spoken in Kaimana Regency, West Papua, Indonesia. Published figures are about 1,000 speakers (a 1991 count cited by Wikipedia) and about 1,500 people (Joshua Project). The searches carried out for this report found only general descriptive and classificatory pages. No primer, textbook, dictionary, grammar, orthography, scripture, audio or school programme for Semimi turned up.

\section{Primary Educational Resources}
None found. The language is not known to be taught in schools, and the searches found no classroom materials.

\section{Linguistic and Reference Materials}
\begin{itemize}
\item Wikipedia, Semimi language (overview, speaker figures): \url{https://en.wikipedia.org/wiki/Semimi_language}
\item Wikipedia, Mairasi languages (family classification, lexical similarity): \url{https://en.wikipedia.org/wiki/Mairasi_languages}
\item New Guinea World, Semimi page: \url{https://newguineaworld.linguistik.uzh.ch/families/etna-bay/semimi}
\item New Guinea World, Etna Bay family page: \url{https://newguineaworld.linguistik.uzh.ch/families/etna-bay/start}
\item ASJP word lists, Semimi 1: \url{https://asjp.clld.org/languages/SEMIMI_1}
\item ASJP word lists, Semimi 2: \url{https://asjp.clld.org/languages/SEMIMI_2}
\end{itemize}

\section{Religious and Community Resources}
Joshua Project lists the Semimi people and gives the writing status as unknown: \url{https://joshuaproject.net/languages/etz}. No scripture translation or community organisation was found.

\section{Digital and Audio Resources}
None found beyond the ASJP word lists above. The Rosetta Project wiki has a stub page for the code: \url{https://rosettapanglossia.longnow.org/index.php?title=Etz}

\section{Curriculum Development Context}
Little is documented. Any curriculum would first need basic documentation work (an orthography, a lexicon, texts recorded with the community), because no teaching or reference material was found. Whether the community wants such a curriculum can only be established by asking the community.
\end{document}